% Generated draft: requires mathematical and source review.
\documentclass[11pt]{article}
\usepackage[margin=25mm]{geometry}
\usepackage{fontspec}
\setmainfont{Latin Modern Roman}
\usepackage{amsmath,amssymb,mathtools,unicode-math}
\setmathfont{Latin Modern Math}
\usepackage{enumitem,needspace,xurl,hyperref}
\setlength{\parindent}{0pt}
\setlength{\parskip}{5pt}
\setlength{\emergencystretch}{3em}
\newcommand{\N}{\mathbb{N}}
\newcommand{\Z}{\mathbb{Z}}
\newcommand{\Q}{\mathbb{Q}}
\newcommand{\R}{\mathbb{R}}
\newcommand{\C}{\mathbb{C}}
\newcommand{\F}{\mathbb{F}}
\newcommand{\A}{\mathbb{A}}
\newcommand{\PP}{\mathbb{P}}
\newcommand{\E}{\mathbb{E}}
\newcommand{\eps}{\varepsilon}
\newcommand{\dd}{\,\mathrm{d}}
\begin{document}
\section*{UnsolvedMath ID: 30001497}
\textbf{Source problem code:} catalogue entry
\subsection{Definitions and mathematical statement}
The supplied question asks whether there exist a polynomial ring \(S\) and a homogeneous ideal \(I\subseteq S\) such that the Hilbert polynomial of \(S/I\) is
\[
p_{S/I}(t)=t^2.
\]
The coefficient ring, grading, and number of variables are not specified in the supplied statement.

Under the conventional interpretation of a polynomial ring with its standard grading over a field, let \(k\) be a field and put \(S=k[x_0,\ldots,x_n]\), where \(n\in\Z_{\geq 0}\) and \(\deg x_j=1\). An ideal \(I\subseteq S\) is homogeneous if it is generated by homogeneous polynomials. Write \((S/I)_d\) for the degree-\(d\) component and define the Hilbert function by
\[
H_{S/I}(d)=\dim_k (S/I)_d,\qquad d\in\Z_{\geq 0}.
\]
The Hilbert polynomial \(p_{S/I}(t)\in\Q[t]\) is the unique polynomial satisfying \(p_{S/I}(d)=H_{S/I}(d)\) for all sufficiently large integers \(d\).

For a specified field \(k\), the corresponding precise existence question is whether
\[
\begin{gathered}
\exists n\in\Z_{\geq 0}\;\exists I\subseteq k[x_0,\ldots,x_n]\text{ homogeneous}\;\exists d_0\in\Z_{\geq 0}\;\forall d\in\Z_{\geq d_0}:\\
\dim_k\bigl(k[x_0,\ldots,x_n]/I\bigr)_d=d^2.
\end{gathered}
\]
This equality is required only in sufficiently large degrees. The source does not specify whether \(k\) is fixed, existentially quantified, or universally quantified, or whether it has additional hypotheses. Consequently, the displayed formulation is conditional on the stated grading and coefficient-field convention; these conventions require confirmation. The accompanying assessment discusses a broader numerical characterization where lex-segment analogues fail, but supplies no formulation of that broader target or its relationship to this existence question.
\subsection{Short English statement}
Ask whether a homogeneous quotient of a polynomial ring can have Hilbert polynomial \(t^2\), meaning that its graded dimensions eventually equal the squares of the degrees under the standard field-based interpretation. The coefficient field, grading conventions, and relationship to the broader background discussion remain unresolved.
\subsection{Sources}
\begin{itemize}
\item Supplied local UnsolvedMath registry record, permanent ID \(30001497\), title \emph{Hilbert Polynomials of Homogeneous Ideals}, exact source problem code \texttt{catalogue entry}; collection: Oberwolfach Reports Open Problems. The record lists proposed year 2010 and names no proposer.
\item Supplied report citation: \emph{Classical Algebraic Geometry} (2010), Oberwolfach Reports, DOI 10.4171/owr/2010/27, pp.\ 1573--1623. \url{https://doi.org/10.4171/owr/2010/27}. No problem-specific page or contributor is supplied.
\item Supplied literature source, labeled ``original OWR Report 27/2010'': \url{https://doi.org/10.4171/owr/2010/27}. The literature assessment was checked on \texttt{2026-08-21}; this is supplied dated curation, not independently verified here. It labels the problem open and describes an unresolved numerical-characterization problem in settings where lex-segment analogues fail, asserting that no general characterization is established. The supplied background explicitly requests independent verification of time-sensitive status and cited claims.
\item The registry assigns statement status \texttt{corrected\_verified} and reports that the clean statement was checked against the official Oberwolfach report and preserves its hypotheses, notation, and open target. These are supplied verification claims, not an independent source check or mathematical review of this draft.
\end{itemize}
\end{document}
